\color{unidarkgreen}{\subsection{Защита от неполнофазного режима}}
\color{black}

\begin{enumerate}

\item
Защита предназначена для ликвидации неполнофазного режима работы выключателей с пофазным приводом, возникающего в результате возможного отказа переключения одной или двух фаз выключателя при коммутации. 

\item
Ввод в работу ЗНР осуществляется с помощью уставки \textbf{«Работа ЗНР»}. Оперативный ввод/вывод ЗНР  производится с помощью БУ (блока управления функциями) \textbf{«Вывод ЗНР »}. Пуск ЗНР формируется при условии появления конфигурируемых сигналов \textbf{«ЗНФ В1»} или \textbf{«ЗНФ В2»} и отключенном положении смежного выключателя (для схем с двумя выключателями), с контролем протекания тока нулевой или обратной последовательности. Выбор количества контролируемых выключателей задается уставкой \textbf{«ЗНР контр. В2»}, \textbf{«ЗНР контр. В3»}. Порог срабатывания ЗНР по току нулевой последовательности задается уставкой \textbf{«Iср ЗНР»}. Выдержка времени срабатывания ЗНР задается уставкой \textbf{«Тср ЗНР »}. 

\item 
Уставки ЗНР приведены в таблице \ref{app_set:znr}. Алгоритм работы ЗНР приведен на рисунке \ref{pict:ZNR}.

\medskip
\begin{figure}[H]
\centering
\includegraphics[width=1\textwidth,height=1\textheight,keepaspectratio]{\fbpath/09.01 БСК/271.0221 ЗНР/img/ZNR.pdf}
\captionof{figure}{Функциональная схема логики ЗНР}\label{pict:ZNR}
\end{figure}


\end{enumerate} 